\documentclass[11pt]{article}
\usepackage[margin=1in]{geometry}
\usepackage{amsmath, amssymb}
\usepackage{enumitem}
\usepackage{hyperref}
\usepackage{fancyhdr}
\usepackage{tcolorbox}

\pagestyle{fancy}
\fancyhf{}
\lhead{Gumbel-Softmax Study Guide}
\rhead{Worksheet 1 of 5}
\cfoot{\thepage}

\newcommand{\E}{\mathbb{E}}
\newcommand{\bz}{\mathbf{z}}
\newcommand{\btheta}{\boldsymbol{\theta}}

\newtcolorbox{keyidea}{colback=blue!5!white, colframe=blue!50!black, title=Key Idea}
\newtcolorbox{definition}{colback=green!5!white, colframe=green!50!black, title=Definition}

\title{\textbf{Worksheet 1: The Problem} \\ \large Why Can't We Backpropagate Through Discrete Variables?}
\author{Based on: \textit{Categorical Reparameterization with Gumbel-Softmax} \\ Jang, Gu, and Poole (2017)}
\date{}

\begin{document}
\maketitle
\thispagestyle{fancy}

\begin{keyidea}
Neural networks are trained via gradient descent, which requires operations in the computation graph to be differentiable. Sampling from a categorical distribution is not differentiable, creating a fundamental obstacle for models with discrete latent variables.
\end{keyidea}

\section*{Learning Objectives}

By the end of this worksheet, you should be able to:

\begin{enumerate}
    \item Explain why categorical sampling blocks ordinary backpropagation.
    \item Distinguish the continuous reparameterization trick from the discrete case.
    \item Compare REINFORCE and straight-through estimators as imperfect workarounds.
    \item State the problem that Gumbel-Softmax is designed to solve.
\end{enumerate}

\section{Background: Categorical Distributions}

\begin{definition}
A categorical distribution over $k$ classes is parameterized by probabilities $\pi_1, \pi_2, \ldots, \pi_k$, where $\pi_i \geq 0$ and $\sum_{i=1}^k \pi_i = 1$. A sample $z$ is often represented as a one-hot vector $\bz \in \{0,1\}^k$ with $\sum_i z_i = 1$.
\end{definition}

\subsection*{Problem 1.1: Warm-Up}

Suppose we have three classes with probabilities $\pi = (0.5,\; 0.3,\; 0.2)$.

\begin{enumerate}[label=(\alph*)]
    \item Write down all possible one-hot samples $\bz$ and their probabilities.
    \item What is $\E[\bz]$?
    \item Is the sample space continuous or discrete? Why does that matter for gradient-based learning?
\end{enumerate}

\section{The Differentiability Problem}

Consider a neural network with a categorical latent variable:

\[
\btheta \longrightarrow \pi = \operatorname{softmax}(f_{\btheta}(x)) \longrightarrow \bz \sim \operatorname{Cat}(\pi) \longrightarrow L = h(\bz)
\]

We want to compute:

\[
\nabla_{\btheta} \E_{\bz \sim \operatorname{Cat}(\pi)}[L(\bz)]
\]

\subsection*{Problem 1.2: Where Does the Gradient Break?}

\begin{enumerate}[label=(\alph*)]
    \item Is $\pi = \operatorname{softmax}(f_{\btheta}(x))$ differentiable with respect to $\btheta$? Explain.
    \item Is the sampling step $\bz \sim \operatorname{Cat}(\pi)$ differentiable with respect to $\pi$? Explain intuitively.
    \item Compare $\pi = (0.49,\; 0.51)$ and $\pi = (0.48,\; 0.52)$. How can the sample change? What does this reveal about trying to assign a gradient to sampling?
\end{enumerate}

\section{The Reparameterization Trick}

\begin{keyidea}
For continuous distributions, the reparameterization trick moves randomness into a fixed noise variable. For example:
\[
\epsilon \sim \mathcal{N}(0,1), \qquad z = \mu + \sigma \epsilon
\]
Now $z$ is a deterministic differentiable function of $\mu$ and $\sigma$ once $\epsilon$ is sampled.
\end{keyidea}

\subsection*{Problem 1.3: Continuous Versus Discrete}

\begin{enumerate}[label=(\alph*)]
    \item In $z = \mu + \sigma\epsilon$, where does the randomness live? Where do the learnable parameters live?
    \item Compute $\frac{\partial}{\partial \mu}\E[z^2]$ using $z = \mu + \sigma\epsilon$ and $\epsilon \sim \mathcal{N}(0,1)$.
    \item Why does this trick not directly solve categorical sampling? What goes wrong if the final operation is $\arg\max$?
\end{enumerate}

\section{Existing Workarounds}

\subsection*{Problem 1.4: REINFORCE}

The score-function estimator uses:

\[
\nabla_{\btheta} \E_{\bz}[f(\bz)] = \E_{\bz}\left[f(\bz)\nabla_{\btheta}\log p_{\btheta}(\bz)\right]
\]

\begin{enumerate}[label=(\alph*)]
    \item Why does REINFORCE avoid differentiating through the sample itself?
    \item Why can this estimator have high variance?
    \item In what kind of model would high-variance gradient estimates become especially painful?
\end{enumerate}

\subsection*{Problem 1.5: Straight-Through Estimation}

The straight-through estimator uses a discrete sample in the forward pass, then pretends the operation was differentiable in the backward pass.

\begin{enumerate}[label=(\alph*)]
    \item Why is this estimator biased?
    \item Why might it still work in practice?
    \item What kind of failure mode would you watch for if training depends heavily on this approximation?
\end{enumerate}

\section{Looking Ahead}

\begin{keyidea}
The Gumbel-Softmax paper asks whether we can build a continuous relaxation of categorical sampling that is differentiable, approaches one-hot categorical samples as temperature decreases, and gives lower-variance gradients than score-function estimators.
\end{keyidea}

\subsection*{Problem 1.6: Preview}

Suppose someone suggests replacing $\arg\max$ with a softmax temperature:

\[
y_i = \frac{\exp(a_i / \tau)}{\sum_j \exp(a_j / \tau)}
\]

\begin{enumerate}[label=(\alph*)]
    \item What happens as $\tau \to 0$?
    \item What happens as $\tau \to \infty$?
    \item Why is temperature alone not enough? What additional ingredient is needed to turn this into a sampling method?
\end{enumerate}

\end{document}

